\section{Linear Attention: reescribir la matriz cuadrática como una recurrencia lineal}

Esta sección explica la intuición básica de Linear Attention. Lo «linear» no significa que el modelo se vuelva una función lineal, sino que su complejidad crece de forma casi lineal con la longitud de la secuencia. Lo habitual es eliminar o reformular la Softmax y, mediante transformaciones algebraicas, escribir la atención como una recurrence similar a la de una RNN. En cada paso se mantiene un estado, de modo que procesar una secuencia de longitud \(L\) cuesta aproximadamente \(O(L)\).

\begin{figure}[H]
\centering
\includegraphics[width=0.94\textwidth]{linear-attention-intuition.png}
\caption{Intuición de Linear Attention: la matriz cuadrática se reescribe como una recurrence lineal. Diagrama conceptual propio, elaborado a partir de la conversación entre 00:07:04 y 00:11:19.}
\end{figure}

\begin{knowledgebox}{Lectura de la figura: Linear Attention no consiste en escribir una capa menos}
Softmax Attention construye de forma explícita una matriz de \(L\times L\). Linear Attention intenta comprimir la información histórica en un estado recurrente, de modo que cada paso de generación solo actualiza ese estado. Así se reducen la complejidad y la presión sobre la KV cache, pero hay que rediseñar la capacidad expresiva y la estabilidad del entrenamiento.
\end{knowledgebox}

\subsection{Su lugar dentro del Transformer}

El Transformer se compone principalmente de bloques de Attention y FFN apilados de forma repetida. En los últimos años, la FFN se transformó con MoE; Linear Attention, en cambio, interviene directamente en el módulo de Attention. Forma parte de la investigación sobre arquitecturas de preentrenamiento, a diferencia del trabajo sobre el optimizer, los datos de preentrenamiento o el post-training. Hoy es más común una arquitectura hybrid: algunas capas conservan Full Attention y otras se sustituyen por Linear Attention.

\begin{importantbox}{Motivación de la arquitectura hybrid}
Una Linear Attention pura puede quedarse corta en capacidad expresiva; una Full Attention pura resulta demasiado costosa. La arquitectura hybrid usa unas pocas capas globales como respaldo de la capacidad expresiva y muchas capas lineales para reducir el costo; hoy es el compromiso más práctico.
\end{importantbox}

\subsection{Resumen de la sección}

El núcleo de Linear Attention es reescribir la complejidad y la representación del estado. Puede reducir el costo de inferencia con contextos largos, pero tiene que resolver la capacidad expresiva, la estabilidad del entrenamiento, el paralelismo en hardware y la validación a gran escala.
